\documentclass[11pt]{article}

% --------------------------------------------------
% Packages
% --------------------------------------------------
\usepackage[T1]{fontenc}
\usepackage{lmodern}
\usepackage{geometry}
\usepackage{microtype}
\usepackage{setspace}
\usepackage{amsmath,amssymb,amsthm}
\usepackage{csquotes}
\usepackage{hyperref}
\usepackage[numbers,authoryear]{natbib}

\geometry{margin=1in}
\setstretch{1.15}

% --------------------------------------------------
% Theorem environments
% --------------------------------------------------
\newtheorem{definition}{Definition}
\newtheorem{proposition}{Proposition}
\newtheorem{remark}{Remark}

% --------------------------------------------------
% Macros
% --------------------------------------------------
\newcommand{\Lang}{\mathcal{L}}
\newcommand{\Sym}{\mathcal{S}}
\newcommand{\Act}{\mathcal{A}}
\newcommand{\Hist}{\mathcal{H}}

% --------------------------------------------------
% Title
% --------------------------------------------------
\title{Constraint Before Reference:\\
Operator Semantics, Reification, and Symbolic Conflict}
\author{Flyxion}
\date{\today}


\begin{document}
\maketitle

\begin{abstract}
Persistent symbolic and ideological conflicts are often attributed to deep
metaphysical or factual disagreement. This paper argues that many such conflicts
are more accurately diagnosed as semantic failures arising from a misordering of
priority between constraint and reference. In a wide class of linguistic,
institutional, and computational systems, meaning functions primarily by
constraining admissible futures rather than by denoting objects. When symbols
that operate as constraints are reclassified as referential absolutes, negotiable
procedural disagreements collapse into non-negotiable existential claims.

The paper develops a constraint-first framework grounded in philosophy of
language, operational semantics, and category theory. It formalizes linguistic
operators as functions on event histories that irreversibly restrict future
possibilities, provides a small-step operational semantics with soundness and
completeness results, and lifts this account to a functorial semantics in which
symbols correspond to endofunctors on spaces of admissible futures. Reification
is shown to constitute a semantic type error with predictable structural
consequences, including loss of negotiability and escalation of conflict.

The framework explains the role of semantic compression and drift in long-lived
symbolic systems, characterizes narrative as a constraint-preserving encoding,
and offers precise diagnostics for distinguishing factual disagreement from
operator misclassification. Applications are developed across legal systems,
organizational governance, computational architectures, and multi-agent
coordination, demonstrating how operator-level reframing can expand resolution
space without requiring consensus on beliefs or interpretations.

By restoring category discipline between constraint and reference, the paper
provides a neutral, formally explicit account of meaning that clarifies why many
disputes persist despite shared normative commitments and identifies structural
conditions under which disagreement can remain negotiable rather than absolute.
\end{abstract}


% ==================================================
\section{Introduction}
% ==================================================

Persistent conflicts across legal, political, organizational, and computational
domains are frequently described as deep, principled, or metaphysical in nature.
They are often attributed to incompatible worldviews, irreconcilable values, or
fundamental disagreements about reality itself. Yet in many such cases, closer
inspection reveals substantial agreement on practical norms governing action,
responsibility, and coordination. What remains unresolved is not what should be
done, but how certain symbols, terms, or roles are understood to function within
the system that governs doing.

This paper advances the thesis that a large class of persistent symbolic
conflicts arise from a specific semantic failure: the inversion of priority
between \emph{constraint} and \emph{reference}. Symbols that originally functioned
to regulate admissible actions, identities, or future states are misclassified as
denoting absolute objects or facts. Once this misclassification occurs,
procedural disagreement is transformed into existential opposition, rendering
negotiation structurally impossible.

The central claim is not that reference is unimportant, but that it is secondary.
Constraint is logically and operationally prior. Meaning stabilizes action before
it stabilizes description.

% --------------------------------------------------
\subsection{A Motivating Case: Jurisdictional Conflict}
% --------------------------------------------------

Consider a recurring form of legal conflict involving jurisdiction. Two courts,
agencies, or institutional bodies may dispute which of them has authority over a
given case. Such disputes are often framed as disagreements about facts: where an
event occurred, which law applies, or which entity possesses legitimate power.

However, jurisdiction does not denote a physical object or discoverable property.
It is an operator: a rule that constrains which actions are permitted by whom and
under what conditions. When jurisdiction is treated as a referential fact rather
than a procedural constraint, disputes escalate. Each party seeks to establish
exclusive ownership of an alleged object (“true authority”), rather than to
negotiate the boundaries of permissible action.

In practice, many jurisdictional disputes are resolved not by discovering new
facts, but by redefining procedural constraints: carving out exceptions,
establishing precedence rules, or introducing shared oversight mechanisms. The
conflict dissolves once the operator structure is restored.

This pattern recurs across domains. The case is representative rather than
exceptional.

% --------------------------------------------------
\subsection{Problem Statement}
% --------------------------------------------------

The prevailing semantic frameworks used to analyze disagreement emphasize
truth-conditions, reference, and correspondence. While powerful, these frameworks
are ill-suited to explaining why conflicts persist in contexts where truth-value
disagreement is minimal or irrelevant.

This paper proposes an alternative explanatory axis. It distinguishes symbols
that function primarily as \emph{operators}—constraining admissible futures—from
symbols that function primarily as \emph{referents}—denoting objects or states of
affairs. Conflicts arise when operator symbols are reinterpreted as referential
absolutes.

This misclassification produces a characteristic set of effects. Procedural norms
are naturalized and treated as features of the world rather than as revisable
rules of coordination, the space of negotiation collapses as alternative
interpretations are excluded, and disagreement is reframed as existential rather
than practical, transforming what could be resolved through procedural revision
into a perceived threat to identity or reality itself.

Understanding this mechanism requires tools beyond standard truth-conditional
semantics.

% --------------------------------------------------
\subsection{Relation to Existing Literature}
% --------------------------------------------------

The framework developed here draws on and synthesizes several distinct traditions.

In philosophy of language, the distinction between sense and reference (Frege),
between description and naming (Russell), and between rigid designation and use
(Kripke) establishes that reference alone cannot account for meaning
\cite{frege1892,russell1905,kripke1980}. Donnellan’s distinction between
referential and attributive uses of descriptions further destabilizes a purely
referential account \cite{donnellan1966}.

Speech act theory and performativity shift attention from truth to action.
Austin’s performatives and Searle’s institutional facts demonstrate that many
utterances do not describe the world but change it \cite{austin1962,searle1969}.
However, these accounts stop short of formalizing long-range constraint on future
states.

Inferentialist and pragmatist approaches, especially Brandom’s account of meaning
as normative inferential role, come closer to the present view
\cite{brandom1994}. Yet they lack a fully explicit temporal or event-based
semantics.

In legal and institutional theory, Hart’s distinction between primary and
secondary rules reveals how authority emerges from constraint rather than fact
\cite{hart1961}. More recent work in institutional design and governance confirms
that symbolic stability depends on enforceable procedural limits, not metaphysical
agreement.

Finally, in computer science, operational semantics, game semantics, and
event-sourced systems already treat meaning as execution and constraint rather
than description. These traditions supply the formal backbone of the present
approach.

This paper unifies these strands into a single, explicit framework.

% --------------------------------------------------
\subsection{Methodology}
% --------------------------------------------------

The methodology of this paper is formal, analytic, and comparative.

First, it introduces a precise semantic distinction between reference and
constraint, expressed using operator semantics over event histories.

Second, it develops an operational semantics in which meaning is evaluated by its
effect on admissible futures rather than by truth conditions.

Third, it lifts this operational account into a categorical framework, allowing
structural comparison between symbolic systems via functors and natural
transformations.

Throughout, formal development is interleaved with concrete examples drawn from
law, organizational systems, and computation. These examples are not illustrative
ornaments; they function as empirical anchors for the formal claims.

The approach is intentionally neutral with respect to metaphysical commitments.
The framework does not deny reference, belief, or interpretation. It specifies
their proper semantic location.

% --------------------------------------------------
\subsection{Contributions and Roadmap}
% --------------------------------------------------

This paper advances several interconnected contributions. It formally
distinguishes between operator symbols and referential symbols, and shows how
systematic misclassification between these semantic types gives rise to
persistent and intractable conflict. It introduces an event-historical
operational semantics in which symbolic meaning is constituted by irreversible
constraints on admissible futures rather than by static denotation. It further
develops a functorial semantics that enables structural comparison between
symbolic systems at an abstract level, independent of their surface
representations. In addition, the analysis identifies structural risks inherent
in highly expressive operator systems, particularly with respect to domination,
lock-in, and loss of negotiability. Finally, the framework is shown to apply
across legal, organizational, and computational domains, demonstrating its
generality and practical relevance.

The paper proceeds as follows. Section~2 deepens the philosophical foundations of
the reference/operator distinction. Section~3 develops a formal account of
semantic compression and drift. Section~4 analyzes reification as a category
error and proves its consequences. Section~5 develops narrative as a constraint-
preserving encoding mechanism. Sections~6 and~7 present the operational and
functorial semantics, respectively. Subsequent sections apply the framework,
compare it with alternatives, and discuss methodological limits and future
directions. The conclusion synthesizes the results and outlines practical and
theoretical implications.

% ==================================================
\section{Names, Reference, and Operators}
% ==================================================

This section develops the conceptual foundation for the paper’s central
distinction between \emph{referential} expressions and \emph{operator} expressions.
The distinction is motivated philosophically, formalized semantically, and
illustrated across legal, institutional, and computational domains. The aim is
not to deny the existence or utility of reference, but to establish that many
conflicts arise when expressions whose primary function is to constrain action
are mistakenly treated as denoting objects.

% --------------------------------------------------
\subsection{Classical Accounts of Reference}
% --------------------------------------------------

Modern philosophy of language begins with the problem of reference. Frege’s
distinction between sense (\emph{Sinn}) and reference (\emph{Bedeutung}) establishes
that linguistic meaning cannot be reduced to denotation alone
\cite{frege1892}. A name may refer to an object while differing in cognitive
content, indicating that reference is only one dimension of meaning.

Russell’s theory of descriptions further complicates the picture by showing that
many apparently referential expressions function quantificationally rather than
denotationally \cite{russell1905}. Kripke’s account of rigid designation restores
a strong notion of reference for proper names, but explicitly limits its scope
\cite{kripke1980}. Even on Kripke’s account, not all grammatically name-like
expressions are rigid designators.

Donnellan’s distinction between referential and attributive uses of descriptions
demonstrates that reference is often a matter of use rather than form
\cite{donnellan1966}. The same expression may function referentially in one
context and non-referentially in another.

These developments collectively undermine the assumption that naming implies
objecthood.

% --------------------------------------------------
\subsection{From Reference to Operation}
% --------------------------------------------------

While the classical literature focuses on reference, a parallel tradition shifts
attention to what expressions \emph{do}. Austin’s theory of speech acts
demonstrates that many utterances perform actions rather than describe states of
affairs \cite{austin1962}. Searle extends this insight to institutional facts,
showing how rules create roles, permissions, and obligations
\cite{searle1969}.

However, speech act theory typically treats performative force as localized to
utterance events. What remains underdeveloped is an account of how symbolic
expressions impose long-range constraints on future action.

To address this gap, we introduce the concept of a linguistic operator.

\begin{definition}[Linguistic Operator]
A linguistic operator is a symbolic expression whose primary semantic function is
to constrain the space of admissible actions, interpretations, or identities
across a structured practice, rather than to denote an object or state of affairs.
\end{definition}

Operators may be realized linguistically, procedurally, or formally. Their
meaning lies in their effect on permissible continuation, not in correspondence
to an external entity.

% --------------------------------------------------
\subsection{Examples Across Domains}
% --------------------------------------------------

The operator/reference distinction appears across many domains.

\paragraph{Legal Domain.}
Terms such as \emph{jurisdiction}, \emph{authority}, and \emph{standing} do not
denote physical objects. They function as operators that delimit who may act, in
what capacity, and under which conditions. Treating such terms as referential
properties invites absolutist conflict; treating them as procedural constraints
enables negotiation and reform \cite{hart1961}.

\paragraph{Institutional Roles.}
Expressions like \emph{citizenship} or \emph{membership} regulate access to rights
and obligations. Their semantic force derives from rule systems, not from
intrinsic properties of individuals. Reifying these roles as natural kinds
produces exclusionary pathologies.

\paragraph{Computational Systems.}
In programming languages, constructs such as \texttt{lock}, \texttt{scope}, or
\texttt{permission} are operators over execution paths. They constrain possible
program states without denoting runtime objects. Misunderstanding these
constructs as values rather than control structures leads to design errors.

These examples illustrate that operator semantics is neither marginal nor exotic.
It is pervasive.

% --------------------------------------------------
\subsection{A Taxonomy of Operator Types}
% --------------------------------------------------

Operators may be classified according to the type of constraint they impose on
admissible continuation.

\begin{definition}[Operator Taxonomy]
Let $\Hist$ denote a space of admissible event histories. Operators may be
distinguished by the manner in which they constrain this space. Some operators
function by excluding subsets of $\Hist$, thereby prohibiting certain future
trajectories. Others operate by authorizing actions that would otherwise be
excluded under existing constraints. A further class imposes precedence
relations, introducing ordering constraints among events without eliminating
them outright. Identification operators act by equating or merging histories,
roles, or trajectories that were previously distinct, while forgetting operators
erase internal distinctions among histories while preserving admissibility at a
coarser level of resolution.
\end{definition}

Each class of operator constrains continuation in a structurally distinct way.
Crucially, none of these operators function by denoting objects; their semantic
force lies entirely in how they reshape the space of admissible futures.


% --------------------------------------------------
\subsection{Criteria for Classification}
% --------------------------------------------------

To determine whether a symbol functions primarily as an operator or as a
referential expression, the framework proposes a set of diagnostic criteria.
A symbol exhibits operator-like behavior when modifications to its use alter
what may occur next without changing the factual description of the present
state, when its interpretation depends essentially on institutional or
procedural context, and when disputes concerning its application can be resolved
through rule revision rather than by appeal to empirical discovery. A further
indicator is temporal reach: symbols that primarily constrain future states of
affairs rather than describe present ones are characteristic of operator
semantics.

Symbols satisfying most of these criteria should therefore be analyzed as
operators rather than as referents.

% --------------------------------------------------
\subsection{Borderline and Hybrid Cases}
% --------------------------------------------------

Some expressions exhibit mixed behavior. For example, \emph{property} may refer
to an asset while simultaneously functioning as an operator regulating use and
transfer. Hybrid cases require careful semantic typing to avoid category error.

\begin{remark}
Hybrid symbols are especially prone to reification errors because their
referential aspect masks their operator function.
\end{remark}

% --------------------------------------------------
\subsection{Historical Development of the Distinction}
% --------------------------------------------------

The gradual recognition of operator semantics parallels broader shifts in
analytic philosophy away from static correspondence models toward use- and
practice-based accounts of meaning. Wittgenstein’s later work emphasizes rule-
following and language games \cite{wittgenstein1953}, while Brandom formalizes
meaning in terms of normative inferential roles \cite{brandom1994}.

The present framework extends these insights by introducing an explicit temporal
and structural account of constraint, enabling precise analysis of persistence,
drift, and conflict.

% --------------------------------------------------
\subsection{Section Summary}
% --------------------------------------------------

This section has argued that many grammatically name-like expressions function
primarily as operators rather than referents. Misclassifying such expressions
produces structural conflict by converting procedural constraints into alleged
objects of dispute. Establishing the operator/reference distinction is therefore
a prerequisite for the formal analysis that follows.

% ==================================================
\section{Semantic Compression and Drift}
% ==================================================

This section develops a rigorous account of semantic compression and drift. The
central claim is that symbolic systems preserve coordination by compressing
high-dimensional normative and procedural constraints into low-dimensional
representations. While such compression is necessary for scalability and
stability, it introduces predictable failure modes. Over time, compressed
symbols drift away from their original operator structure, increasing the risk
of reification and conflict.

The analysis proceeds in four stages. First, semantic compression is formalized
using information-theoretic tools. Second, historical case studies illustrate
long-term drift in institutional symbols. Third, a typology of drift mechanisms
is developed. Finally, semantic drift is modeled as a dynamical process over
constraint spaces.

% --------------------------------------------------
\subsection{Compression as an Information-Theoretic Necessity}
% --------------------------------------------------

Let a symbolic system be understood as a mapping from complex constraint
structures to manageable representations. Coordination across agents requires
that these representations be communicable, memorable, and repeatable.

Formally, let $X$ denote a space of high-dimensional constraint states, and let
$Y$ denote a space of symbolic representations. A semantic compression is a
mapping
\[
C : X \to Y
\]
that reduces descriptive complexity while preserving sufficient information for
coordination.

Using Shannon entropy, the compression ratio may be characterized as
\[
\kappa = \frac{H(Y)}{H(X)},
\]
where $H(\cdot)$ denotes entropy. For long-lived symbolic systems, $\kappa \ll 1$
is not optional but necessary.

However, compression is lossy. Information about the original constraint
topology is discarded. The meaning of a symbol is therefore underdetermined by
its form alone.

\begin{definition}[Semantic Compression]
Semantic compression is the lossy encoding of a high-dimensional constraint
structure into a low-dimensional symbolic representation that preserves
coordination-relevant information while discarding contextual detail.
\end{definition}

This definition applies equally to legal terms, organizational roles, technical
abstractions, and computational interfaces.

% --------------------------------------------------
\subsection{Decompression and Interpretation}
% --------------------------------------------------

Symbols are not used directly. They are interpreted, or decompressed, back into
constraints on action.

Let
\[
D_t : Y \to X
\]
be a decompression function at time $t$. Crucially, $D_t$ is not fixed. It
depends on background practices, incentives, and institutional context.

Meaning at time $t$ is therefore given by the composite
\[
D_t \circ C.
\]

Semantic drift occurs when $D_t$ changes while $C$ remains fixed. Because the
compressed symbol does not encode its original context, reinterpretation is
unavoidable.

% --------------------------------------------------
\subsection{Historical Case Studies of Drift}
% --------------------------------------------------

\paragraph{Constitutional Terminology.}
Terms such as \emph{due process} or \emph{equal protection} originated as
procedural constraints on governance. Over centuries, they have accumulated
layers of interpretation, judicial precedent, and political pressure. Disputes
over their meaning are rarely resolved by discovering facts; they are resolved by
reconfiguring operator structure through reinterpretation and amendment.

\paragraph{Scientific Concepts.}
Scientific terms such as \emph{gene}, \emph{information}, or \emph{fitness} have
undergone substantial drift as empirical knowledge and theoretical frameworks
evolved. Early compressed meanings proved insufficient, leading to contested
reinterpretations. Treating these terms as fixed referents rather than evolving
operators has repeatedly generated confusion.

\paragraph{Organizational Roles.}
Roles such as \emph{manager} or \emph{fiduciary} encode complex responsibility
structures. As organizations scale and technologies change, these compressed
symbols drift away from their original constraint sets, producing governance
failures unless explicitly re-articulated.

These cases demonstrate that drift is not accidental; it is structural.

% --------------------------------------------------
\subsection{Mechanisms of Semantic Drift}
% --------------------------------------------------

Semantic drift arises through multiple interacting mechanisms that alter how
compressed symbolic structures are interpreted over time.

\begin{definition}[Drift Mechanisms]
Semantic drift refers to systematic changes in the interpretation of symbolic
operators induced by shifts in the conditions of their use. Common mechanisms
include the loss of background practices that originally anchored interpretation,
reinterpretation driven by misaligned incentives or strategic advantage, the
imposition of preferred decompressions by agents or institutions with asymmetric
power, changes in feasible action brought about by new technological
affordances, and the gradual accumulation of institutional modifications that
layer partial adjustments rather than re-articulating underlying constraints.
\end{definition}


Each mechanism alters $D_t$ without modifying the symbol itself, increasing
semantic divergence across agents.

% --------------------------------------------------
\subsection{Drift as a Dynamical System}
% --------------------------------------------------

Let $D_t$ evolve according to a dynamical process
\[
D_{t+1} = F(D_t, I_t, P_t),
\]
where $I_t$ represents incentive structures and $P_t$ represents power
distributions.

The resulting system may exhibit a range of dynamical behaviors over time,
including convergence toward a stable interpretation, oscillation corresponding
to persistently contested meaning, divergence in which interpretations fragment
across subcommunities, and lock-in characterized by path-dependent rigidity. In
this dynamical perspective, reification corresponds to a particular attractor in
which an operator is treated as a fixed referent, thereby freezing the evolution
of the admissible future space $D_t$ and eliminating the possibility of further
negotiation.


\begin{remark}
Reification is not semantic stability but semantic arrest.
\end{remark}

% --------------------------------------------------
\subsection{Compression--Drift Tradeoff}
% --------------------------------------------------

There exists a fundamental tradeoff between compression efficiency and drift
risk. Highly compressed symbols scale well but drift rapidly; richly articulated
symbols resist drift but scale poorly.

This tradeoff can be expressed schematically as:
\[
\text{Scalability} \;\uparrow \quad \Longleftrightarrow \quad \text{Drift Risk}
\;\uparrow
\]

Managing symbolic systems therefore requires periodic re-articulation of
operator structure.

% --------------------------------------------------
\subsection{Section Summary}
% --------------------------------------------------

Semantic compression is necessary for coordination but introduces unavoidable
drift. Drift becomes pathological when compressed operator symbols are treated
as referential absolutes, producing reification and conflict. Understanding drift
as a dynamical process prepares the ground for the formal analysis of category
error in the next section.

% ==================================================
\section{Reification and Category Error}
% ==================================================

This section analyzes reification as a precise semantic failure: a category error
in which linguistic operators are misclassified as referential expressions. The
analysis proceeds by formal definition, proof of consequences, examination of
boundary cases, and articulation of strategies for reversing reification once it
has occurred.

The central result is that reification transforms negotiable constraint systems
into absolutist symbolic regimes by collapsing procedural disagreement into
existential opposition.

% --------------------------------------------------
\subsection{Formal Definition of Reification}
% --------------------------------------------------

Let $\Lang$ be a language whose expressions are typed into distinct semantic
categories. We distinguish in particular between two primary types. Expressions
of type $\mathcal{R}$ are referential, and their semantic evaluation proceeds by
correspondence to objects, properties, or states of affairs. Expressions of type
$\mathcal{O}$ are operators, and their semantic contribution consists in
constraining the space of admissible future continuations of an event history
rather than in denoting entities.


\begin{definition}[Semantic Typing]
A semantic typing is a function
\[
\tau : \Lang \to \{\mathcal{R}, \mathcal{O}\}
\]
assigning each expression its dominant semantic role.
\end{definition}

\begin{definition}[Reification Error]
A reification error occurs when an expression $\sigma$ such that
$\tau(\sigma) = \mathcal{O}$ is treated, interpreted, or adjudicated as if
$\tau(\sigma) = \mathcal{R}$.
\end{definition}

Reification is thus not a metaphysical mistake but a semantic one.

% --------------------------------------------------
\subsection{Formal Consequences of Reification}
% --------------------------------------------------

We now prove the central proposition introduced earlier.

\begin{proposition}[Reification Produces Absolutism]
If a symbolic system treats its highest-order operators as referential objects,
then disagreement within that system tends toward absolutism rather than
negotiation.
\end{proposition}

\begin{proof}
Let $\sigma$ be an operator constraining admissible histories
$\sigma : \Hist \to \mathcal{P}(\Hist)$. As an operator, disagreement over $\sigma$
concerns which constraints to impose and is resolvable by modification,
composition, or reinterpretation.

If $\sigma$ is retyped as referential, disagreement over $\sigma$ becomes a
dispute over the existence, identity, or nature of an object. Referential
disagreement admits no procedural resolution; it can only be settled by denial
or elimination of opposing claims.

Thus, the space of negotiable futures collapses. Disagreement becomes existential
rather than procedural. \qed
\end{proof}

% --------------------------------------------------
\subsection{Corollaries}
% --------------------------------------------------

\begin{corollary}[Collapse of Negotiation Space]
Reification strictly reduces the set of admissible resolutions available to a
symbolic system.
\end{corollary}

\begin{corollary}[Path Dependence]
Once reification occurs, subsequent reinterpretations accumulate asymmetrically,
producing lock-in effects.
\end{corollary}

\begin{corollary}[Conflict Amplification]
Reified operators amplify conflict by converting rule disagreement into identity
threat.
\end{corollary}

Each corollary follows directly from the elimination of procedural degrees of
freedom.

% --------------------------------------------------
\subsection{Boundary Conditions and Counterexamples}
% --------------------------------------------------

Not all symbols are susceptible to reification in the same way.

\paragraph{Pure Referents.}
Expressions denoting physical objects or well-defined quantities are correctly
typed as referential. Treating them as operators produces confusion but not
absolutism.

\paragraph{Explicitly Typed Operators.}
Systems that explicitly distinguish rules from facts—such as typed programming
languages or well-designed legal codes—exhibit strong resistance to reification.

\paragraph{Hybrid Symbols.}
Symbols combining referential and operator aspects require careful handling.
Failure to disambiguate their roles produces instability rather than immediate
absolutism.

These cases delimit the scope of the theory.

% --------------------------------------------------
\subsection{Cognitive and Social Mechanisms of Reification}
% --------------------------------------------------

Reification is reinforced by well-documented cognitive and social tendencies.

\paragraph{Cognitive Load Reduction.}
Treating operators as objects simplifies reasoning by reducing temporal and
procedural complexity.

\paragraph{Essentialist Bias.}
Humans exhibit a bias toward treating abstract roles as natural kinds, even when
they are rule-defined.

\paragraph{Institutional Incentives.}
Power asymmetries reward actors who present procedural constraints as immutable
facts.

These mechanisms explain why reification is common even when formally unjustified.

% --------------------------------------------------
\subsection{De-Reification: Recovering Operator Structure}
% --------------------------------------------------

Despite its stabilizing appearance, reification is reversible.

\begin{definition}[De-Reification]
De-reification is the semantic operation of restoring an expression’s operator
typing by re-articulating the constraints it imposes on admissible futures.
\end{definition}

Effective strategies for de-reification involve making governing rules explicit,
reframing disputes in procedural rather than referential terms, and decomposing
constraints temporally so that their effects on admissible futures can be
examined over time. In addition, reconstructing the original operator intent
through narrative analysis can restore the regulative function of symbols that
have become reified, thereby reopening interpretive and negotiable space.


Crucially, de-reification does not require consensus on interpretation—only
agreement that the symbol functions procedurally.

% --------------------------------------------------
\subsection{Formal Characterization of De-Reification}
% --------------------------------------------------

Let $\sigma$ be reified. De-reification corresponds to replacing a referential
judgment
\[
\sigma = x
\]
with an operator judgment
\[
\sigma(h) = \{ h' \mid h' \text{ satisfies constraint } C \}.
\]

This restores negotiability by reopening the space of admissible histories.

\begin{remark}
De-reification reintroduces time into meaning. Absolutism is timeless; procedure
is temporal.
\end{remark}

% --------------------------------------------------
\subsection{Section Summary}
% --------------------------------------------------

Reification is a semantic category error with predictable structural
consequences. By mistaking operators for referents, symbolic systems collapse
negotiation into absolutism. Recognizing and reversing this error is a necessary
precondition for durable coordination. The next section examines how narrative
serves as a constraint-preserving alternative to reification.

% ==================================================
\section{Narrative as Constraint-Preserving Encoding}
% ==================================================

The preceding sections analyzed the failure modes that arise when symbolic
operators are reified as referents. This section develops the complementary
thesis: narrative functions as a robust mechanism for preserving operator
structure without collapsing it into absolutist form. Narrative encodes
constraints indirectly, distributing semantic force across time, agents, and
consequence rather than locating it in isolated propositions.

This section formalizes narrative as a constraint-preserving encoding over event
histories, compares it to alternative preservation mechanisms, and explains its
central role in maintaining negotiability in long-lived symbolic systems.

% --------------------------------------------------
\subsection{Narrative Versus Propositional Encoding}
% --------------------------------------------------

Propositional encoding represents meaning by asserting facts. Narrative encoding
represents meaning by exhibiting sequences of events subject to constraints.

Let $\Hist$ denote a space of event histories. A propositional statement
partitions $\Hist$ into histories where the proposition holds and histories where
it does not. By contrast, a narrative selects a structured subset of $\Hist$
whose admissible continuations implicitly encode norms.

\begin{definition}[Narrative Encoding]
A narrative encoding of a constraint $C$ is a structured projection
\[
N \subseteq \Hist
\]
such that admissible interpretations of $N$ enforce $C$ without explicitly
asserting it as a proposition.
\end{definition}

Narratives thus preserve constraint topology while avoiding explicit operator
invocation.

% --------------------------------------------------
\subsection{Temporal Structure and Constraint Distribution}
% --------------------------------------------------

Narratives distribute meaning across temporal structure. No single event or
utterance carries the full semantic load. Instead, constraints emerge from
patterns of action and consequence.

This distribution has two important effects. First, it resists reification by
preventing meaning from being localized in any single symbol, assertion, or
moment. Second, it permits multiple decompressions that remain consistent with
the same underlying operator structure, allowing interpretive flexibility
without sacrificing semantic coherence.

Formally, let $N = \langle e_1, \dots, e_n \rangle$ be a narrative sequence. The
constraint enforced by $N$ is not attributable to any $e_i$ alone but to the
relations among them.

\begin{remark}
Narrative meaning is irreducibly temporal. Attempts to extract a timeless core
invite reification.
\end{remark}

% --------------------------------------------------
\subsection{Narrative and Temporal Logic}
% --------------------------------------------------

Narrative constraints can be formalized using temporal logic. Let $\Box$ denote
necessity over futures and $\Diamond$ possibility.

Rather than asserting $\Box P$ directly, a narrative may demonstrate that any
history in which $P$ fails leads to undesirable or excluded continuations. The
constraint is encoded by exhibited consequence rather than asserted necessity.

This aligns narrative semantics with model-checking rather than truth
evaluation.

% --------------------------------------------------
\subsection{Comparison with Other Preservation Mechanisms}
% --------------------------------------------------

Narrative is not the only method for preserving operator structure, but it has
distinct advantages.

\paragraph{Ritualized Procedure.}
Procedures encode constraints through repetition but are brittle under context
change.

\paragraph{Case Law and Precedent.}
Precedent preserves constraint through exemplars, closely approximating narrative
structure. Its success depends on analogical reasoning.

\paragraph{Formal Specification.}
Formal systems preserve constraints precisely but scale poorly across human
interpretive communities.

Narrative occupies a unique position: high interpretive flexibility with
constraint stability.

% --------------------------------------------------
\subsection{Examples Across Domains}
% --------------------------------------------------

\paragraph{Legal Reasoning.}
Judicial opinions often rely on narrative reconstruction of facts to justify
procedural outcomes. The binding force lies not in explicit rule citation but in
the exhibited pattern of reasoning.

\paragraph{Organizational Memory.}
Organizations transmit norms through stories of success and failure. These
stories encode constraints on acceptable behavior without issuing formal
directives.

\paragraph{Simulation and Training.}
Simulations encode constraints by allowing agents to explore outcome spaces.
Narrative scenarios function as executable operator semantics.

% --------------------------------------------------
\subsection{Narrative and Reinterpretation}
% --------------------------------------------------

Because narrative avoids explicit operator invocation, it supports reinterpretive
plasticity. Different agents may extract different surface meanings while
preserving deep constraint structure.

Let $D_i$ and $D_j$ be two decompression functions applied to the same narrative
$N$. If both yield constraint-equivalent interpretations, coordination remains
possible despite disagreement.

\begin{definition}[Constraint-Equivalent Interpretations]
Two interpretations are constraint-equivalent if they induce the same admissible
future space.
\end{definition}

This explains how narrative sustains coordination across heterogeneous agents.

% --------------------------------------------------
\subsection{Computational Perspectives}
% --------------------------------------------------

From a computational standpoint, narrative corresponds to execution traces in
which constraints are inferred from the distinction between reachable and
unreachable states rather than asserted propositionally. On this view, narrative
semantics aligns naturally with formalisms such as process algebras, game
semantics, reinforcement learning environments, and scenario-based planning,
each of which characterizes meaning in terms of permitted trajectories through a
state space. In all such cases, semantic content is not given by isolated
assertions but emerges from the structure of possible and impossible paths.

% --------------------------------------------------
\subsection{Narrative as Anti-Reification Mechanism}
% --------------------------------------------------

Narrative resists reification by distributing meaning across temporal structure,
by embedding constraints in patterns of consequence rather than in isolated
assertions, and by preventing individual symbols from being detached from the
contexts that give them regulative force. Attempts to extract a single
authoritative proposition from a narrative typically destroy its semantic
function, replacing constraint-guided interpretation with a brittle and
misleading referential reading.

\begin{remark}
Narrative fails only when it is forced to answer referential questions it was not
designed to answer.
\end{remark}

% --------------------------------------------------
\subsection{Section Summary}
% --------------------------------------------------

Narrative functions as a constraint-preserving encoding that avoids the category
errors associated with reification. By distributing semantic force across time
and consequence, it stabilizes operator structure while maintaining
negotiability. This makes narrative indispensable to long-lived symbolic systems.
The next sections formalize these insights in operational and categorical terms.

% ==================================================
\section{Operational Semantics on Event Histories}
% ==================================================

This section provides a complete operational semantics for symbolic meaning
understood as constraint on admissible futures. Rather than assigning truth
conditions to expressions, the semantics specifies how symbolic operators act on
event histories to restrict or transform the space of possible continuations.

Operational semantics is chosen for two reasons. First, it makes temporal
irreversibility explicit. Second, it aligns semantic evaluation with execution,
rendering meaning directly actionable.

% --------------------------------------------------
\subsection{Event Histories}
% --------------------------------------------------

Let $\Hist$ be a set of finite, prefix-closed sequences of events. An event
history $h \in \Hist$ is written as
\[
h = \langle e_1, e_2, \dots, e_n \rangle,
\]
where each $e_i$ represents an irreversible update.

\begin{definition}[Prefix Order]
For histories $h, h' \in \Hist$, we write $h \preceq h'$ if $h$ is a prefix of $h'$.
\end{definition}

The prefix order induces a partial order on $\Hist$ representing irreversible
extension.

\begin{definition}[Future Space]
For a history $h$, the future space is
\[
\mathsf{Fut}(h) = \{ h' \in \Hist \mid h \preceq h' \}.
\]
\end{definition}

Once an event occurs, histories incompatible with it are permanently excluded
from $\mathsf{Fut}(h)$.

% --------------------------------------------------
\subsection{Operators as Transition Constraints}
% --------------------------------------------------

An operator acts by eliminating or restructuring futures rather than by adding
content to history.

\begin{definition}[Operator]
An operator is a function
\[
\sigma : \Hist \to \mathcal{P}(\Hist)
\]
such that $\sigma(h) \subseteq \mathsf{Fut}(h)$ for all $h \in \Hist$.
\end{definition}

Operators are monotone: they cannot introduce futures incompatible with the
current history.

\begin{lemma}[Monotonicity]
For any operator $\sigma$ and histories $h \preceq h'$,
\[
\sigma(h') \subseteq \sigma(h).
\]
\end{lemma}

\begin{proof}
Any future admissible after $h'$ must also be admissible after $h$, since
$h'$ introduces additional constraints. \qed
\end{proof}

% --------------------------------------------------
\subsection{Configurations and Rewrite Rules}
% --------------------------------------------------

Semantic evaluation proceeds by rewriting configurations.

\begin{definition}[Configuration]
A configuration is a pair $(h, \Sigma)$ where $h \in \Hist$ and $\Sigma$ is a
multiset of active operators.
\end{definition}

A rewrite relation
\[
(h, \Sigma) \longrightarrow (h', \Sigma')
\]
specifies a single semantic step.

% --------------------------------------------------
\subsection{Core Rewrite Schemes}
% --------------------------------------------------

We define representative rewrite schemes corresponding to common operator types.

\paragraph{POP (Exclusion).}
\[
\frac{h' \in \mathsf{Fut}(h) \quad h' \notin \sigma(h)}
{(h, \{\sigma\}) \longrightarrow (h, \emptyset)}
\]

POP eliminates incompatible futures.

\paragraph{BIND (Precedence).}
\[
(h, \{\textsc{Bind}(e_i \prec e_j)\}) \longrightarrow (h, \Sigma')
\]
where $\Sigma'$ enforces the ordering constraint on all extensions of $h$.

\paragraph{MERGE (Identification).}
\[
(h_1, h_2, \{\textsc{Merge}\}) \longrightarrow (h, \Sigma)
\]
where $h$ is the least upper bound of $h_1$ and $h_2$ if one exists.

\paragraph{COLLAPSE (Forgetting).}
\[
(h, \{\textsc{Collapse}\}) \longrightarrow (h', \Sigma)
\]
where $h'$ is a quotient of $h$ under an equivalence relation preserving
admissibility.

Each rule eliminates or restructures futures without adding new ones.

% --------------------------------------------------
\subsection{Soundness}
% --------------------------------------------------

\begin{theorem}[Soundness]
If $(h, \Sigma) \longrightarrow (h', \Sigma')$, then
\[
\mathsf{Fut}(h') \subseteq \mathsf{Fut}(h).
\]
\end{theorem}

\begin{proof}
Each rewrite rule either preserves $h$ or replaces it with a quotient or merge
consistent with $h$. No rule introduces histories incompatible with prior events.
Therefore admissible futures are monotonically reduced. \qed
\end{proof}

Soundness ensures that semantic evaluation corresponds to constraint
strengthening.

% --------------------------------------------------
\subsection{Completeness}
% --------------------------------------------------

Completeness ensures that all constraint reductions can be realized operationally.

\begin{theorem}[Completeness]
For any operator $\sigma$ and history $h$, if $h' \in \sigma(h)$, then there exists
a finite rewrite sequence
\[
(h, \{\sigma\}) \longrightarrow^{*} (h', \Sigma')
\]
realizing the constraint imposed by $\sigma$.
\end{theorem}

\begin{proof}
By construction, operators specify admissible futures. Each elimination or
restructuring step corresponds to a rewrite rule. Finite composition suffices to
exclude all incompatible histories. \qed
\end{proof}

Completeness guarantees that operator semantics is executable.

% --------------------------------------------------
\subsection{Worked Examples}
% --------------------------------------------------

Consider a history $h = \langle e_1 \rangle$ and an operator enforcing that event
$e_2$ may occur only after $e_3$.

Applying \textsc{Bind} yields:
\[
(h, \{\textsc{Bind}(e_3 \prec e_2)\}) \longrightarrow (h, \Sigma')
\]
where $\Sigma'$ excludes all futures in which $e_2$ precedes $e_3$.

The resulting future space contains only histories consistent with the
constraint. No propositional assertion is required.

% --------------------------------------------------
\subsection{Comparison with Other Semantic Frameworks}
% --------------------------------------------------

\paragraph{Denotational Semantics.}
Denotational approaches assign mathematical objects to expressions. While
powerful, they obscure temporal irreversibility and negotiation dynamics.

\paragraph{Modal Semantics.}
Modal logic captures necessity and possibility but treats modality statically.
Operator semantics embeds modality in execution.

\paragraph{Game Semantics.}
Game semantics models interaction explicitly. Operator semantics may be viewed as
a restriction of game semantics to admissible trajectories.

Operational semantics uniquely combines temporal structure with constraint
evaluation.



% --------------------------------------------------
\subsection{Algebraic Properties of Operators}
% --------------------------------------------------

Operators compose naturally.

\begin{definition}[Operator Composition]
Given operators $\sigma_1, \sigma_2$, define
\[
(\sigma_2 \circ \sigma_1)(h) = \sigma_2(\sigma_1(h)).
\]
\end{definition}

Composition is associative but not necessarily commutative. Identity operators
act as neutral elements.

These properties prepare the ground for categorical abstraction.

% --------------------------------------------------
\subsection{Section Summary}
% --------------------------------------------------

This section established a complete operational semantics for meaning as
constraint on event histories. Soundness and completeness ensure that semantic
evaluation is both faithful and executable. The next section lifts this framework
to a categorical level, enabling structural comparison between symbolic systems.

% ==================================================
\section{From Operator Semantics to Event--Historical Structure}
% ==================================================

This appendix provides a precise correspondence between the operator-based
semantic framework developed above and an event-historical semantics in which
meaning, identity, and obligation are constituted by irreversible transitions.

% ==================================================
\subsection{Event--Historical Substrate}
% ==================================================

Let $\Hist$ denote a partially ordered set of events equipped with a precedence
relation $\prec$. An event history $h \in \Hist$ is a finite or countably infinite
prefix-closed sequence of events.

\begin{definition}[Event]
An event is an irreversible update that reduces the space of admissible future
histories.
\end{definition}

Irreversibility is essential: an event does not merely describe change but
\emph{settles} what can no longer occur.

% ==================================================
\subsection{Operators as Constraints on Futures}
% ==================================================

Recall that an operator was defined as a function constraining admissible action.
We now refine this definition in event-historical terms.

\begin{definition}[Event-Historical Operator]
An operator $\sigma$ is a function
\[
\sigma : \Hist \to \mathcal{P}(\Hist)
\]
such that for any history $h$, $\sigma(h)$ is the set of admissible future
extensions of $h$ consistent with $\sigma$.
\end{definition}

Operators do not add content to history; they restrict continuation. Their
semantic force is entirely future-directed.

This formalizes the intuition that sacred or legal symbols regulate conduct rather
than denote entities.

% ==================================================
\subsection{Spherepop Correspondence}
% ==================================================

In Spherepop-style calculi, existence and identity arise from explicit operations
rather than axioms. We may map core Spherepop primitives to the present framework
as follows.

\begin{center}
\begin{tabular}{ll}
\textbf{Spherepop Operation} & \textbf{Semantic Interpretation} \\
\midrule
POP & Irreversible elimination of future branches \\
MERGE & Identification of histories under shared constraints \\
BIND & Introduction of precedence constraints on events \\
COLLAPSE & Erasure of historical differentiation \\
\end{tabular}
\end{center}

Each operation corresponds to a transformation on $\Hist$ that alters the space
of admissible continuations.

\begin{proposition}
Spherepop operations are higher-order operators acting on event histories rather
than on values or propositions.
\end{proposition}

\begin{proof}
Each operation modifies $\mathcal{P}(\Hist)$ by restricting, identifying, or
releasing future extensions. No operation asserts a truth condition; all impose
structural constraints.
\end{proof}

% ==================================================
\subsection{Identity as Stabilized History}
% ==================================================

Under this mapping, identity is not primitive.

\begin{definition}[Historical Identity]
An identity is an equivalence class of event histories stabilized under a set of
operators.
\end{definition}

Names, titles, and symbolic markers do not refer to entities but to stabilized
regions of $\Hist$. Reification occurs when such equivalence classes are mistaken
for timeless objects.

This explains why identity-based conflicts intensify when symbolic operators are
treated as ontological predicates.


% ==================================================
\section{A Functorial Semantics of Constrained Futures}
% ==================================================

The previous section defined meaning operationally as constraint on admissible
event histories. This section abstracts that account categorically. The resulting
functorial semantics makes explicit which structural features of symbolic systems
are invariant under reinterpretation and which are not.

Category theory is employed not as an added layer of abstraction, but as a means
of isolating semantic structure from representational accident. Two symbolic
systems are semantically equivalent, on this view, if their constraint structures
are categorically isomorphic, regardless of surface vocabulary.

% --------------------------------------------------
\subsection{The Category of Event Histories}
% --------------------------------------------------

Define a category $\mathbf{Hist}$ whose objects are event histories
$h \in \Hist$ and whose morphisms encode irreversible extension. For any two
histories $h$ and $h'$, there exists a unique morphism $f : h \to h'$ if and only
if $h \preceq h'$, that is, if $h$ is a prefix of $h'$.

\begin{proposition}
The category $\mathbf{Hist}$ is thin, and hence a preorder, encoding the
irreversible extension of event histories.
\end{proposition}


\begin{proof}
By definition, there is at most one morphism between any two objects. Reflexivity
and transitivity follow from prefix inclusion. \qed
\end{proof}

The categorical structure directly mirrors temporal irreversibility.

% --------------------------------------------------
\subsection{The Category of Constraint Spaces}
% --------------------------------------------------

Let $\mathbf{Constr}$ be the category defined by:

\begin{itemize}
\item Objects are subsets $C \subseteq \Hist$ closed under extension
(i.e.\ if $h \in C$ and $h \preceq h'$, then $h' \in C$).
\item Morphisms are monotone functions preserving inclusion.
\end{itemize}

\begin{proposition}
$\mathbf{Constr}$ is a complete lattice category.
\end{proposition}

\begin{proof}
Arbitrary intersections and unions of extension-closed sets are extension-closed,
yielding limits and colimits. \qed
\end{proof}

$\mathbf{Constr}$ captures admissibility structure independently of presentation.

% --------------------------------------------------
\subsection{The Meaning Functor}
% --------------------------------------------------

We now define the central semantic construction.

\begin{definition}[Meaning Functor]
Define a functor
\[
\mathcal{M} : \mathbf{Hist} \to \mathbf{Constr}
\]
by
\[
\mathcal{M}(h) = \mathsf{Fut}(h),
\]
and for $f : h \to h'$,
\[
\mathcal{M}(f)(C) = C \cap \mathsf{Fut}(h').
\]
\end{definition}

\begin{theorem}[Functoriality]
$\mathcal{M}$ is a well-defined functor.
\end{theorem}

\begin{proof}
Identity morphisms map to identity inclusions. Composition of prefix morphisms
corresponds to successive restriction of admissible futures. Associativity
follows from set intersection. \qed
\end{proof}

The meaning functor assigns to each history its space of admissible continuations.

% --------------------------------------------------
\subsection{Operators as Endofunctors}
% --------------------------------------------------

Each operator induces an endofunctor on $\mathbf{Constr}$.

\begin{definition}[Operator Endofunctor]
Given an operator $\sigma$, define an endofunctor
\[
F_\sigma : \mathbf{Constr} \to \mathbf{Constr}
\]
by
\[
F_\sigma(C) = C \cap \sigma(h),
\]
for any history $h$ generating $C$.
\end{definition}

\begin{proposition}
Operator composition corresponds to functor composition.
\end{proposition}

\begin{proof}
Intersection of admissible futures is associative. Identity operators act as
identity functors. \qed
\end{proof}

Thus, symbolic systems correspond to monoids of endofunctors.

% --------------------------------------------------
\subsection{Natural Transformations as Reinterpretation}
% --------------------------------------------------

Reinterpretation of symbols corresponds categorically to natural transformation.

\begin{definition}[Semantic Reinterpretation]
A reinterpretation of an operator $\sigma$ as $\sigma'$ is a natural
transformation
\[
\eta : F_\sigma \Rightarrow F_{\sigma'}.
\]
\end{definition}

Naturality ensures that reinterpretation commutes with history extension.

\begin{remark}
Semantic drift corresponds to loss of naturality: reinterpretation behaves
differently across contexts.
\end{remark}

This formalizes the intuition that drift is not mere change, but structural
misalignment.

% --------------------------------------------------
\subsection{Limits, Colimits, and Identity}
% --------------------------------------------------

MERGE operations correspond to colimits in $\mathbf{Constr}$, while COLLAPSE
corresponds to coequalizers.

\begin{proposition}
Stable identities arise as fixed points of operator endofunctor composition.
\end{proposition}

\begin{proof}
An identity persists exactly when repeated application of constraints leaves the
admissible future space invariant. \qed
\end{proof}

Identity is thus an emergent categorical invariant, not a primitive referent.

% --------------------------------------------------
\subsection{Semantic Equivalence}
% --------------------------------------------------

Two symbolic systems are semantically equivalent if their meaning functors are
naturally isomorphic.

\begin{definition}[Semantic Equivalence]
Two systems $S_1$ and $S_2$ are semantically equivalent if there exists a natural
isomorphism
\[
\mathcal{M}_1 \cong \mathcal{M}_2.
\]
\end{definition}

This criterion ignores vocabulary and focuses solely on constraint structure.

\begin{remark}
Many apparent metaphysical disagreements reduce to failures of recognizing
semantic equivalence.
\end{remark}

% --------------------------------------------------
\subsection{Adjunctions and Relaxation}
% --------------------------------------------------

Relaxation of constraints corresponds to left adjoints; enforcement corresponds
to right adjoints.

\begin{definition}[Constraint Adjunction]
An adjunction $L \dashv R$ between endofunctors represents reversible negotiation
between permissiveness and restriction.
\end{definition}

Adjunctions formalize the possibility of controlled loosening without collapse.

% --------------------------------------------------
\subsection{Section Summary}
% --------------------------------------------------

This section lifted operational semantics into a categorical framework. Meaning is
a functor from histories to admissible futures; symbols are endofunctors; and
reinterpretation is natural transformation. This abstraction isolates semantic
invariants and clarifies when disagreement is substantive and when it is merely
presentational.

The next section applies the framework to concrete domains and extended case
studies.

% ==================================================
\section{Applications and Case Studies}
% ==================================================

The preceding sections developed a formal theory in which meaning is understood as
constraint on admissible futures rather than reference to objects. This section
demonstrates the explanatory and diagnostic power of that framework across
multiple domains. In each case, persistent conflict or instability is shown to
arise from reification of operator symbols, and resolution is shown to depend on
restoring operator-level interpretation.

The aim is not to offer policy prescriptions, but to show how semantic structure
conditions the space of possible resolutions.

% --------------------------------------------------
\subsection{Legal Systems}
% --------------------------------------------------

Legal systems provide a paradigmatic case of operator semantics at scale. Legal
language is replete with expressions that regulate conduct without denoting
objects.

\paragraph{Constitutional Interpretation.}
Constitutional terms such as \emph{authority}, \emph{due process}, or
\emph{jurisdiction} function as operators that constrain permissible actions of
institutions. Disputes framed as disagreements about the “true meaning” of such
terms often persist because the terms are treated as referential absolutes rather
than adjustable constraints.

Under an operator semantics, constitutional interpretation becomes a problem of
constraint preservation under changing conditions. Amendments, precedent, and
doctrinal tests function as reinterpretations—natural transformations—rather than
discoveries of fixed reference.

\paragraph{Contract Law.}
Contracts encode future constraints on action. Treating contractual language as
descriptive rather than performative leads to brittle enforcement and adversarial
interpretation. Operator semantics clarifies why doctrines such as good faith and
reasonableness are indispensable: they reintroduce negotiability at the operator
level.

% --------------------------------------------------
\subsection{Organizational Governance}
% --------------------------------------------------

Organizations coordinate action through role definitions, procedures, and
protocols. These are operator-rich systems.

\paragraph{Role Reification.}
Roles such as \emph{manager}, \emph{executive}, or \emph{stakeholder} regulate
decision authority. When these roles are reified as intrinsic personal properties
rather than procedural positions, organizations experience rigidity, blame
cycles, and power struggles.

Operator semantics explains why effective organizations continuously revise role
definitions: they are re-articulating constraints to maintain adaptability.

\paragraph{Protocol Design.}
Organizational protocols function as executable narratives. Case escalation
procedures, incident response playbooks, and governance charters encode
constraints via conditional sequences rather than propositions.

Failures often arise when protocols are treated as literal instructions rather
than as constraint-generating structures adaptable to context.

% --------------------------------------------------
\subsection{Computational Systems}
% --------------------------------------------------

Computational systems provide a domain where operator semantics is already
implicitly accepted.

\paragraph{Programming Language Semantics.}
Control structures such as scopes, locks, and permissions constrain execution
paths. Treating these constructs as values rather than operators leads to security
and concurrency failures.

Operational semantics makes explicit why static typing and access control are
forms of semantic constraint enforcement.

\paragraph{Distributed Systems.}
Consensus protocols do not describe truth; they restrict admissible histories of
system state. Operator semantics clarifies why such systems fail under network
partition: admissible future spaces fragment.

\paragraph{AI and Multi-Agent Systems.}
Reward functions and policies are operators over agent trajectories. Reifying
objectives as fixed goals rather than adjustable constraints produces brittle and
misaligned systems. Operator semantics reframes alignment as constraint
negotiation over futures rather than goal specification.

% --------------------------------------------------
\subsection{Conflict Resolution and Mediation}
% --------------------------------------------------

Many entrenched conflicts persist despite agreement on outcomes because parties
disagree on symbolic framing.

Operator semantics suggests a diagnostic procedure for addressing persistent
symbolic conflict. The procedure begins by identifying symbols that are treated as
referential absolutes rather than as regulative devices. It then proceeds by
recovering the operator functions those symbols originally served, followed by
the explicit re-articulation of the constraints they impose on admissible action.
Finally, negotiability is restored through procedural revision, allowing
coordination to occur without requiring agreement on underlying metaphysical or
interpretive commitments.

This procedure does not require consensus on narratives or identities. It requires
only agreement that certain symbols regulate action rather than denote objects.

% --------------------------------------------------
\subsection{Formal Modeling of a Case Study}
% --------------------------------------------------

Let two agents share a history $h$ and disagree over a symbol $\sigma$. If
$\sigma$ is treated referentially, admissible futures partition into disjoint
sets with no overlap.

Under operator interpretation, both agents agree on a constraint
$\sigma : \Hist \to \mathcal{P}(\Hist)$, even if they differ on surface
interpretation. The intersection
\[
\mathsf{Fut}_1(h) \cap \mathsf{Fut}_2(h)
\]
remains non-empty, preserving coordination.

This formal distinction explains why reframing disputes procedurally often
succeeds where factual argument fails.

% --------------------------------------------------
\subsection{Section Summary}
% --------------------------------------------------

Across legal, organizational, and computational systems, persistent conflict
arises when operator symbols are reified. Resolution depends not on discovering
correct reference, but on restoring constraint structure. The next section
compares this framework with alternative semantic theories and addresses
potential objections.

% ==================================================
\section{Comparison with Alternative Frameworks}
% ==================================================

This section situates the constraint-first framework relative to established
approaches in philosophy of language, semantics, and coordination theory. The
goal is not to dismiss alternative frameworks, but to clarify the specific
phenomena each explains well, the failure modes they exhibit, and the sense in
which operator semantics provides a complementary—and in some cases strictly
stronger—explanatory account of persistent symbolic conflict.

% --------------------------------------------------
\subsection{Truth-Conditional Semantics}
% --------------------------------------------------

Truth-conditional semantics evaluates meaning by specifying the conditions under
which an expression is true or false. This approach has been extraordinarily
successful in modeling descriptive language and factual discourse.

However, truth-conditional semantics encounters difficulty in domains where
disagreement persists despite agreement on observable facts. In such cases,
truth conditions are either trivially satisfied or irrelevant to the dispute.

Operator semantics explains this failure mode: the disputed expressions do not
primarily function to describe states of affairs. They regulate admissible
actions and identities. Evaluating them by truth conditions misclassifies their
semantic role.

Truth-conditional semantics remains appropriate for referential expressions. The
present framework argues only that it is incomplete as a general theory of
meaning.

% --------------------------------------------------
\subsection{Pragmatist and Inferentialist Accounts}
% --------------------------------------------------

Pragmatist approaches emphasize use, practice, and consequences over reference.
Inferentialist theories, particularly those developed by Brandom, understand
meaning in terms of normative inferential roles.

These accounts align closely with the present framework. Both reject the
primacy of correspondence and emphasize rule-governed practice. However, they
typically lack an explicit temporal or event-historical semantics.

Constraint-first semantics extends inferentialist approaches by making the role of
irreversibility explicit, by modeling how constraints operate over extended
temporal horizons rather than only locally within inferential contexts, and by
providing an executable semantics in which meaning can be realized through
concrete transitions over event histories. The result is a more precise account
of how meaning stabilizes over time and of why semantic misclassification, rather
than substantive disagreement, so often gives rise to persistent conflict.

% --------------------------------------------------
\subsection{Speech Act Theory and Performativity}
% --------------------------------------------------

Speech act theory demonstrates that many utterances perform actions rather than
describe facts. Performativity captures how language brings about institutional
states.

While foundational, speech act theory remains largely local: it focuses on the
effect of utterances at the moment of performance. Operator semantics generalizes
this insight to persistent, distributed constraint over histories.

In effect, operator semantics treats institutions themselves as extended speech
acts whose force unfolds temporally.

% --------------------------------------------------
\subsection{Game-Theoretic and Economic Models}
% --------------------------------------------------

Game-theoretic models analyze coordination in terms of incentives, equilibria,
and strategies. These models excel at predicting behavior under fixed rules.

However, they typically assume the rule structure as given. When disputes arise
over the rules themselves, game-theoretic analysis must step outside its formal
frame.

Operator semantics addresses this meta-level by modeling how rules constrain
admissible futures and how those constraints may be renegotiated. It complements
game theory by explaining how the game itself is constituted and transformed.

% --------------------------------------------------
\subsection{Realist Objections}
% --------------------------------------------------

A common objection holds that some symbolic disputes concern real facts and
cannot be resolved procedurally. The present framework does not deny this.
Rather, it provides a criterion for distinguishing factual disagreement from
semantic misclassification.

Where disagreement turns on empirical claims, reference remains primary. Where
disagreement persists despite empirical convergence, operator semantics offers a
diagnostic alternative.

Constraint-first semantics thus narrows the scope of realist claims without
rejecting realism outright.

% --------------------------------------------------
\subsection{Summary of Comparative Advantages}
% --------------------------------------------------

The distinctive contributions of the present framework lie in its explicit
semantic typing of operators as distinct from referential expressions, its
event-historical treatment of meaning as constraint on admissible futures, and
its formal modeling of reification and de-reification as structurally identifiable
semantic transformations. In addition, it provides categorical criteria for
semantic equivalence grounded in preserved constraint structure and offers
practical diagnostic tools for analyzing and mitigating persistent symbolic
conflict. No existing framework integrates all of these elements within a single,
coherent formal account.

% --------------------------------------------------
\subsection{Section Summary}
% --------------------------------------------------

Constraint-first semantics does not replace truth-conditional, pragmatist, or
game-theoretic approaches. It explains a class of failures they systematically
leave unresolved. The next section addresses methodological considerations,
scope, and limitations of the framework.

% ==================================================
\section{Methodological Considerations}
% ==================================================

This section addresses methodological questions raised by a constraint-first
approach to meaning. It clarifies what neutrality entails in this context,
identifies the limits of formal modeling, specifies the scope of applicability,
and discusses how the framework can be empirically evaluated. These
considerations are essential for situating the theory within broader research
programs and for avoiding category errors in its application.

% --------------------------------------------------
\subsection{Neutrality and Non-Commitment}
% --------------------------------------------------

The framework developed in this paper is neutral with respect to metaphysical,
ideological, and interpretive commitments. Neutrality here does not mean
indifference to outcomes, nor does it imply value skepticism. It means that the
theory does not presuppose any particular account of what exists, what is true,
or what ultimately matters.

Instead, the framework specifies how symbolic expressions function within
rule-governed systems to constrain admissible futures. This functional focus
allows the theory to diagnose semantic failures without adjudicating underlying
beliefs.

\begin{remark}
Neutrality is achieved by semantic typing, not by suspension of judgment.
\end{remark}

By distinguishing operator semantics from referential semantics, the framework
avoids implicit metaphysical commitments while preserving analytical precision.

% --------------------------------------------------
\subsection{Limits of Formal Modeling}
% --------------------------------------------------

Formal semantics necessarily abstracts away from aspects of lived practice.
While the operational and functorial models capture structural features of
constraint, they do not encode affect, persuasion, or power directly.

This limitation is deliberate. The framework is designed to analyze the
\emph{conditions of negotiability}, not to model every factor influencing
decision-making. Psychological, sociological, and political variables enter the
analysis indirectly, through their effect on interpretation functions and power
asymmetries.

Formal clarity requires selective abstraction.

% --------------------------------------------------
\subsection{Scope of Applicability}
% --------------------------------------------------

The constraint-first framework applies in contexts where symbols function
primarily to regulate action rather than to describe objects, where coordination
depends on shared procedural constraints rather than on shared representations,
and where disagreement persists even after substantial convergence on empirical
facts. In such settings, semantic failure is best diagnosed at the level of
constraint interpretation rather than referential disagreement.


These conditions obtain in legal systems, organizational governance,
computational protocols, and multi-agent environments. The framework is less
useful for domains dominated by empirical discovery or purely descriptive
language.

\begin{remark}
Applying operator semantics to genuinely referential disputes constitutes a
category error.
\end{remark}

The theory therefore includes its own criteria for appropriate use.

% --------------------------------------------------
\subsection{Generalizability and Transfer}
% --------------------------------------------------

Generalizability is achieved structurally rather than substantively. The same
formal apparatus applies across domains because it abstracts from domain-specific
content and models constraint topology instead.

Transfer between domains requires careful retyping of symbols and explicit
identification of admissible histories. Failure to perform this retyping risks
importing inappropriate assumptions.

% --------------------------------------------------
\subsection{Empirical Testability}
% --------------------------------------------------

Although the framework is formal, it yields empirically testable predictions.

First, it predicts that conflicts framed in referential terms will exhibit
intractability even when procedural consensus exists. Second, it predicts that
reframing disputes at the operator level will expand the space of negotiated
outcomes.

These predictions can be tested through a combination of experimental studies
that manipulate negotiation framing, comparative analyses of institutional
reforms across differing semantic regimes, and simulations of multi-agent systems
operating under alternative constraint structures. Each of these approaches
provides a means of assessing how semantic typing and operator reinterpretation
affect the persistence or resolution of conflict.

Empirical validation thus concerns patterns of resolution and persistence, not
truth-apt judgments.

% --------------------------------------------------
\subsection{Falsifiability}
% --------------------------------------------------

The framework is falsifiable in several ways. If conflicts diagnosed as
reification failures remain intractable after operator-level reframing, the
theory is challenged. If referential reframing consistently outperforms
procedural reframing in operator-dominated domains, the central claim is
undermined.

The theory therefore makes substantive, testable claims about the structure of
disagreement.

% --------------------------------------------------
\subsection{Methodological Summary}
% --------------------------------------------------

Constraint-first semantics is a structural, neutral, and formally explicit
approach to meaning. Its power lies in diagnosing semantic failure modes and
clarifying the conditions under which negotiation is possible. Its limits are
those of abstraction itself.

The next section outlines future research directions opened by this framework.

% ==================================================
\section{Future Directions}
% ==================================================

The constraint-first framework developed in this paper opens multiple avenues for
further research. Some involve deepening the formal apparatus, others involve
empirical validation, and still others involve practical deployment in
computational and institutional systems. This section outlines these directions
and identifies open questions that remain unresolved.

The aim is not to forecast outcomes, but to clarify where the theory’s internal
logic points next.

% --------------------------------------------------
\subsection{Computational Implementations}
% --------------------------------------------------

One immediate direction is the construction of executable systems that embody
operator semantics directly. Event-sourced architectures, rewrite systems, and
policy engines already approximate this structure, but often without explicit
semantic typing.

Future work may include the development of programming languages that treat
operators as first-class semantic entities, the design of runtime systems that
explicitly track and manipulate spaces of admissible futures, and the creation
of formal verification tools capable of reasoning about the preservation and
compatibility of constraints over execution histories. In addition, interpreter
designs that make processes of de-reification explicit as program
transformations represent a promising direction for rendering semantic structure
transparent within computational systems.

Such systems would make semantic structure visible and manipulable, reducing the
risk of unintentional reification in software design.

% --------------------------------------------------
\subsection{Multi-Agent Systems and Coordination}
% --------------------------------------------------

The framework naturally extends to multi-agent environments in which coordination
depends on shared constraints rather than on shared beliefs or representations.
Within such settings, several open questions arise concerning the dynamics of
interaction and interpretation. These include how agents negotiate
reinterpretations of governing operators, how spaces of admissible futures
fragment or merge under conditions of partial agreement, how asymmetries of power
shape the distribution and enforcement of constraints, and how stable patterns of
coordination emerge through iterative processes of de-reification and
reinterpretation.

These questions can be explored through simulation, game-theoretic extension,
and experimental economics.

% --------------------------------------------------
\subsection{Empirical Research Programs}
% --------------------------------------------------

Empirical validation of the framework necessarily requires interdisciplinary
collaboration, drawing on methods from the social sciences, cognitive science,
and computational analysis. Promising lines of inquiry include controlled
experimental studies that manipulate negotiation framing, longitudinal analyses
of institutional reform processes, corpus-based investigations of symbolic drift
in legal or organizational texts, and behavioral studies examining susceptibility
to reification bias in interpretive judgment.

Across these approaches, the central empirical variable of interest is not
convergence on truth-apt claims, but the expansion or contraction of the space of
negotiable futures made available to participants under different semantic
framings.

% --------------------------------------------------
\subsection{Formal Extensions}
% --------------------------------------------------

Several formal questions remain open.

\paragraph{Complexity.}
What is the computational complexity of determining non-empty intersection of
admissible future spaces under multiple operators?

\paragraph{Decidability.}
Under what conditions is constraint compatibility decidable?

\paragraph{Higher-Order Operators.}
What are the properties of systems in which operators act on other operators?

\paragraph{Topological Structure.}
Can constraint spaces be given a meaningful topology or metric capturing degrees
of negotiability?

Addressing these questions would strengthen the mathematical foundations of the
theory.

% --------------------------------------------------
\subsection{Institutional Design}
% --------------------------------------------------

The framework suggests a set of design principles for institutions and protocols
concerned with durable coordination. In particular, institutions are more likely
to remain adaptable and less prone to absolutist conflict when they explicitly
distinguish symbols that function as operators from those that function as
referential terms, when they incorporate mechanisms for the periodic
de-reification and reinterpretation of their governing symbols, and when they
maintain a non-empty space of admissible future actions under revision and
change.

The formal articulation and evaluation of these principles constitutes an
important direction for future applied research.

% --------------------------------------------------
\subsection{Integration with Other Theories}
% --------------------------------------------------

Constraint-first semantics is compatible with, but distinct from, a range of
existing theoretical frameworks. Future work may explore deeper integration with
inferentialist approaches to meaning, with models from institutional economics
and governance theory that emphasize rule-structured coordination, and with
formal epistemology concerned with belief revision and informational dynamics.
In addition, the framework invites further engagement with research on AI
alignment, where the management of admissible futures and the avoidance of
semantic reification are increasingly recognized as central design challenges.


Such integration should proceed cautiously, preserving semantic typing
discipline.

% --------------------------------------------------
\subsection{Section Summary}
% --------------------------------------------------

The theory developed here is not complete. It is deliberately foundational. Its
value lies in identifying a semantic failure mode and providing tools to analyze
and correct it. The directions outlined above represent natural extensions rather
than speculative leaps.

The concluding section synthesizes the results and articulates their broader
significance.

% ==================================================
\section{Conclusion}
% ==================================================

This paper has advanced a simple but far-reaching thesis: in a wide class of
symbolic systems, constraint is logically and operationally prior to reference.
Many persistent conflicts arise not from deep disagreement about facts, values,
or reality, but from a semantic misclassification in which operator symbols are
treated as referential absolutes. Once this inversion occurs, negotiable
procedural disagreement collapses into existential opposition.

The framework developed here restores category discipline by distinguishing
symbols that denote objects from symbols that constrain admissible futures. It
does not deny reference, belief, or interpretation. It specifies their proper
place.

% --------------------------------------------------
\subsection{Summary of Results}
% --------------------------------------------------

The paper has established the following results:

First, it developed a principled distinction between referential expressions and
operator expressions, grounded in the philosophy of language and extended through
formal semantics. This distinction explains why many grammatically name-like
symbols resist truth-conditional analysis.

Second, it showed that long-lived symbolic systems rely on semantic compression,
which is necessary for coordination but introduces predictable drift. Drift
becomes pathological when compressed operators are reified as fixed referents.

Third, it provided a formal account of reification as a semantic category error
and proved that reification collapses negotiation space by converting procedural
constraints into alleged objects of dispute.

Fourth, it demonstrated that narrative functions as a constraint-preserving
encoding that resists reification by distributing semantic force across time and
consequence.

Fifth, it introduced an operational semantics on event histories in which meaning
is executable constraint, and proved soundness and completeness of the resulting
rewrite system.

Sixth, it lifted this operational account into a functorial semantics, showing
that symbols correspond to endofunctors, reinterpretation to natural
transformations, and identity to fixed points of constraint composition.

Seventh, it applied the framework across legal, organizational, computational,
and multi-agent domains, demonstrating its diagnostic and explanatory power.

Eighth, it situated the framework relative to existing semantic, pragmatic, and
coordination theories, clarifying both its compatibility and its distinct
contributions.

% --------------------------------------------------
\subsection{Theoretical Contributions}
% --------------------------------------------------

The principal theoretical contributions of this work can be summarized as
follows. The paper introduces a constraint-first ordering of semantic priority
that applies across institutional, computational, and interpretive domains. It
develops a formal account of reification and de-reification as semantic type
errors and corrective reinterpretations, respectively. It advances an
event-historical semantics in which meaning is constituted by irreversible
constraints on admissible futures rather than by static reference. It provides a
categorical criterion for semantic equivalence grounded in preserved constraint
structure, and it offers a structural explanation of why certain symbolic
conflicts persist despite apparent normative convergence.


Together, these contributions provide a unified account of meaning, stability,
and disagreement that neither reduces symbols to reference nor dissolves them
into indeterminacy.

% --------------------------------------------------
\subsection{Practical Implications}
% --------------------------------------------------

The practical implications of a constraint-first semantics are substantial.

For legal and institutional systems, the framework clarifies why procedural
re-articulation often succeeds where factual argument fails. For organizations,
it explains the necessity of revisiting role definitions and protocols. For
computational and AI systems, it reframes alignment as the management of
admissible futures rather than the specification of fixed objectives.

In each case, stability depends on preserving negotiability at the operator
level.

% --------------------------------------------------
\subsection{Ethical Considerations and Misuse}
% --------------------------------------------------

The expressive power of operator semantics entails ethical risk. Systems capable
of constraining futures at scale can be used to stabilize cooperation or to
eliminate it. The difference lies not in the formalism but in how operator
composition is governed.

This paper has therefore emphasized transparency, explicit typing, and
preservation of non-empty admissible future space as minimal structural
safeguards. These are not moral add-ons; they are conditions for avoiding semantic
capture.

% --------------------------------------------------
\subsection{Philosophical Significance}
% --------------------------------------------------

At a broader level, the framework challenges a deeply ingrained assumption in the
theory of meaning: that reference is primary and constraint derivative. The
analysis here suggests the opposite ordering in many of the domains where meaning
matters most for coordination.

Meaning stabilizes action before it stabilizes description. Symbols bind futures
before they describe worlds.

Recognizing this priority does not resolve all disagreement. It does something
more fundamental: it restores the possibility of resolution by returning
conflict to the domain of procedure rather than absolutes.

% --------------------------------------------------
\subsection{Closing Remark}
% --------------------------------------------------

Constraint before reference is not a slogan but a discipline. It is a demand for
semantic precision in the places where imprecision carries the highest cost. By
recovering the operator structure of symbolic systems, we recover the conditions
under which disagreement can remain productive rather than destructive.

That recovery is not optional. It is a precondition for durable coordination in
any system complex enough to endure.

\newpage
% ==================================================
\appendix
\section{Formal System Specification}
% ==================================================

This appendix presents the core formal system underlying the constraint-first
semantics developed in the main body. The purpose of this appendix is threefold:
(1) to make all semantic assumptions explicit, (2) to provide a typed foundation
that prevents category errors, and (3) to prepare the ground for operational and
categorical proofs in subsequent appendices.

% ==================================================
\subsection{Semantic Typing Discipline}
% ==================================================

We begin by fixing the semantic universe and its typing discipline.

\subsubsection{Semantic Types}

\begin{definition}[Semantic Universe]
The semantic universe consists of three disjoint domains:
\[
\mathcal{R} \quad \text{(referential terms)}, \qquad
\mathcal{O} \quad \text{(operators)}, \qquad
\mathcal{H} \quad \text{(event histories)}.
\]
These domains are semantically disjoint. In particular, operators are not
objects, and histories are not propositions.
\end{definition}

\begin{remark}
Semantic typing is independent of grammatical form. A syntactic noun phrase may
denote an element of $\mathcal{O}$ rather than $\mathcal{R}$.
\end{remark}

This separation is the formal mechanism by which reification errors are detected.

% ==================================================
\subsection{Event Histories}
% ==================================================

Event histories encode irreversible structure.

\begin{definition}[Event]
An event $e$ is an atomic, irreversible update that eliminates at least one
possible future continuation.
\end{definition}

\begin{definition}[Event History]
An event history $h$ is a finite or countably infinite sequence
\[
h = \langle e_1, e_2, \dots \rangle
\]
equipped with a precedence relation that is transitive and antisymmetric, and
that satisfies prefix closure. These conditions ensure that the ordering of
events is well-defined, irreversible, and compatible with the interpretation of
histories as settled pasts extended by admissible futures.
\end{definition}


\begin{definition}[Prefix Order]
For histories $h, h' \in \mathcal{H}$, define
\[
h \preceq h' \iff h \text{ is a prefix of } h'.
\]
\end{definition}

Prefix order encodes irreversibility: once an event occurs, all incompatible
histories are excluded.

% ==================================================
\subsection{Admissible Futures}
% ==================================================

\begin{definition}[Future Space]
For a history $h$, define its future space as
\[
\mathsf{Fut}(h) = \{ h' \in \mathcal{H} \mid h \preceq h' \}.
\]
\end{definition}

\begin{lemma}[Monotonicity of Futures]
If $h \preceq h'$, then
\[
\mathsf{Fut}(h') \subseteq \mathsf{Fut}(h).
\]
\end{lemma}

\begin{proof}
Any continuation of $h'$ is necessarily a continuation of $h$, since $h$ is a
prefix of $h'$. No new futures are introduced by extension.
\end{proof}

This lemma formalizes temporal asymmetry.

% ==================================================
\subsection{Operators as Future Constraints}
% ==================================================

\begin{definition}[Operator]
An operator is a function
\[
\sigma : \mathcal{H} \to \mathcal{P}(\mathcal{H})
\]
such that
\[
\sigma(h) \subseteq \mathsf{Fut}(h) \quad \text{for all } h \in \mathcal{H}.
\]
Operators do not generate events. They constrain which future histories remain
admissible.
\end{definition}

\begin{lemma}[Operator Monotonicity]
If $h \preceq h'$, then
\[
\sigma(h') \subseteq \sigma(h).
\]
\end{lemma}

\begin{proof}
Since $\mathsf{Fut}(h') \subseteq \mathsf{Fut}(h)$ and
$\sigma(h') \subseteq \mathsf{Fut}(h')$, it follows that
$\sigma(h') \subseteq \mathsf{Fut}(h)$. Operators cannot re-admit excluded
futures.
\end{proof}

% ==================================================
\subsection{Operator Composition}
% ==================================================

\begin{definition}[Operator Composition]
Given operators $\sigma_1, \sigma_2 \in \mathcal{O}$, define their composition by
\[
(\sigma_2 \circ \sigma_1)(h) = \sigma_2(h) \cap \sigma_1(h).
\]
\end{definition}

\begin{proposition}[Closure]
The set $\mathcal{O}$ is closed under composition.
\end{proposition}

\begin{proof}
Intersection preserves admissibility: both $\sigma_1(h)$ and $\sigma_2(h)$ are
subsets of $\mathsf{Fut}(h)$, hence so is their intersection.
\end{proof}

Composition is commutative but not necessarily idempotent.

% ==================================================
\subsection{Reification as a Type Error}
% ==================================================

\begin{definition}[Reification Error]
A reification error occurs when an expression of semantic type $\mathcal{O}$ is
treated as if it were of type $\mathcal{R}$.
\end{definition}

\begin{lemma}[Semantic Collapse Under Reification]
If an operator $\sigma$ is reified, then its interpretation no longer constrains
$\mathsf{Fut}(h)$ but instead introduces truth-conditional evaluation.
\end{lemma}

\begin{proof}
Truth-conditional evaluation presupposes denotation. Operators lack denotation and
only define admissibility. Reification therefore discards constraint semantics and
replaces it with an inapplicable evaluation regime.
\end{proof}

This formalizes why reification collapses negotiation space: admissibility is
replaced by alleged facticity.

% ==================================================
\subsection{Minimal Consistency Conditions}
% ==================================================

\begin{definition}[Semantic Consistency]
A symbolic system is semantically consistent if and only if for all histories
$h$,
\[
\bigcap_{\sigma \in \Sigma} \sigma(h) \neq \varnothing.
\]
\end{definition}

\begin{remark}
Consistency here is structural, not logical. It concerns future availability,
not contradiction.
\end{remark}

% ==================================================
\subsection{Appendix Summary}
% ==================================================

This appendix has established a typed semantic universe in which referential
expressions, operators, and event histories are formally distinguished. Within
this universe, event histories were characterized as irreversible structures,
and operators were defined as constraints on admissible future continuations
rather than as denotational entities. Reification was identified as a formal type
error arising from the misclassification of operators as referential terms, and
semantic consistency was defined in structural terms as the preservation of a
non-empty future space.

These definitions provide a sufficient foundation for the operational semantics
and accompanying proofs developed in subsequent sections, ensuring that later
results rest on a precise and internally coherent formal basis.

\newpage
% ==================================================
\section{Operational Semantics and Expanded Proofs}
% ==================================================

This appendix develops the operational semantics implicit in the main text and
provides detailed proofs of soundness, completeness, and related properties.
Where the main body emphasized conceptual structure, this appendix makes the
execution model explicit.

% ==================================================
\subsection{Configurations and Execution State}
% ==================================================

\subsubsection{Configurations}

\begin{definition}[Configuration]
A configuration is an ordered pair $(h, \Sigma)$, where $h \in \mathcal{H}$
denotes the current event history and $\Sigma \subseteq \mathcal{O}$ is a finite
multiset of operators active at that history. The configuration represents a
semantic state consisting of a settled past together with the constraints that
currently govern the space of admissible future continuations.
\end{definition}


\begin{remark}
Configurations do not contain values, propositions, or truth assignments. They
encode only history and constraint.
\end{remark}

% ==================================================
\subsection{Small-Step Transition Relation}
% ==================================================

We define a small-step transition relation
\[
(h, \Sigma) \;\longrightarrow\; (h', \Sigma').
\]

Each transition rule satisfies two invariants:
\begin{enumerate}
\item \emph{Prefix preservation}: $h \preceq h'$,
\item \emph{Constraint monotonicity}: admissible futures do not expand.
\end{enumerate}

The transition relation formalizes execution as progressive restriction of
future possibility.

% ==================================================
\subsection{Primitive Rewrite Rules}
% ==================================================

We define the following canonical rule schemas.

% --------------------------------------------------
\subsubsection{POP (Elimination)}
% --------------------------------------------------

The POP rule removes futures violating a constraint.

\[
\frac{
h' \in \mathsf{Fut}(h)
\quad
h' \notin \sigma(h)
}{
(h, \{\sigma\} \cup \Sigma)
\;\longrightarrow\;
(h, \Sigma)
}
\]

This rule enforces the constraint without altering history. It eliminates
incompatible continuations but introduces no new structure.

% --------------------------------------------------
\subsubsection{BIND (Precedence Constraint)}
% --------------------------------------------------

The BIND rule introduces ordering constraints on future events.

\[
(h, \{\textsc{BIND}(e_i \prec e_j)\} \cup \Sigma)
\;\longrightarrow\;
(h, \Sigma')
\]

where $\Sigma'$ contains an operator enforcing the precedence relation
$e_i \prec e_j$ on all future extensions.

BIND does not settle events; it constrains their admissible orderings.

% --------------------------------------------------
\subsubsection{MERGE (Identification)}
% --------------------------------------------------

The MERGE rule identifies compatible histories.

\[
(h_1, h_2, \textsc{MERGE})
\;\longrightarrow\;
(h, \Sigma)
\]

where $h$ is the least upper bound of $h_1$ and $h_2$ under prefix order, if such a
history exists.

MERGE fails when no consistent joint history exists, producing deadlock.

% --------------------------------------------------
\subsubsection{COLLAPSE (Forgetting)}
% --------------------------------------------------

The COLLAPSE rule erases internal distinctions while preserving admissibility.

\[
(h, \textsc{COLLAPSE})
\;\longrightarrow\;
(h', \Sigma)
\]

where $h'$ is a quotient of $h$ under an equivalence relation that preserves
future availability.

COLLAPSE reduces historical resolution without increasing possibility.

% ==================================================
\subsection{Invariant Preservation}
% ==================================================

\begin{lemma}[Prefix Preservation]
For any transition $(h, \Sigma) \rightarrow (h', \Sigma')$, we have
\[
h \preceq h'.
\]
\end{lemma}

\begin{proof}
POP, BIND, and COLLAPSE do not extend history. MERGE produces a history that
contains both input histories as prefixes. No rule removes past events.
\end{proof}

% ==================================================
\subsection{Soundness Theorem}
% ==================================================

\begin{theorem}[Soundness]
If
\[
(h, \Sigma) \longrightarrow (h', \Sigma'),
\]
then
\[
\mathsf{Fut}(h') \subseteq \mathsf{Fut}(h).
\]
\end{theorem}

\begin{proof}
The claim follows by inspection of the rewrite rules. In the case of
\textsc{POP}, admissible futures are restricted by intersection with the
constraint set $\sigma(h)$. The \textsc{BIND} rule excludes precisely those
future histories that violate the imposed precedence relations. The
\textsc{MERGE} rule retains only those histories that are jointly consistent
with the merged inputs, while the \textsc{COLLAPSE} rule preserves admissible
futures up to an explicitly defined equivalence relation. In all cases, the
effect of rewriting is either to preserve or to eliminate elements of
$\mathsf{Fut}(h)$; no rule introduces histories that were not already admissible.
It follows that admissible futures are preserved or strictly reduced under every
transition.
\end{proof}

Soundness therefore guarantees that semantic execution respects the
irreversibility of event histories.


% ==================================================
\subsection{Completeness Theorem}
% ==================================================

\begin{theorem}[Completeness]
Let $C \subseteq \mathsf{Fut}(h)$ be any constraint expressible as a finite
intersection of operator images. Then there exists a finite sequence of rewrite
steps yielding a configuration whose admissible futures are exactly $C$.
\end{theorem}

\begin{proof}
By assumption,
\[
C = \bigcap_{i=1}^n \sigma_i(h).
\]
Applying the corresponding operators via POP in any order yields the desired
restriction. Order is irrelevant due to commutativity and associativity of
intersection.
\end{proof}

Completeness ensures that the operational semantics can realize all expressible
constraints.

% ==================================================
\subsection{Deadlock and Collapse}
% ==================================================

\begin{definition}[Semantic Deadlock]
A configuration $(h, \Sigma)$ is in deadlock if
\[
\bigcap_{\sigma \in \Sigma} \sigma(h) = \varnothing.
\]
\end{definition}

\begin{lemma}[Deadlock Persistence]
Once deadlock is reached, no further transitions are possible.
\end{lemma}

\begin{proof}
All rewrite rules presuppose non-empty admissible futures. When none exist, no
rule applies.
\end{proof}

Deadlock corresponds to semantic absolutism: no continuation is permitted.

% ==================================================
\subsection{Determinacy and Confluence}
% ==================================================

\begin{proposition}[Confluence]
Operator application is confluent: the order of operator enforcement does not
affect the resulting admissible future space.
\end{proposition}

\begin{proof}
Operator composition is set intersection, which is associative and commutative.
All reduction paths therefore converge to the same admissible future space.
\end{proof}

Confluence guarantees determinacy of semantic outcome.

% ==================================================
\subsection{Appendix Summary}
% ==================================================

This appendix has developed a small-step operational semantics for the
constraint-first framework and demonstrated that semantic execution preserves
the irreversibility invariants imposed by event histories. It has established
soundness and completeness results for the execution model, showing that all and
only admissible future restrictions can be realized through finite sequences of
rewrite steps. In addition, the appendix has provided a formal characterization
of semantic deadlock in terms of empty future spaces and shown that operator
application is confluent, ensuring that the order of constraint enforcement does
not affect the resulting admissible futures.

These results justify treating meaning as executable constraint rather than
denotation and provide the technical backbone for the categorical semantics
developed elsewhere.

\newpage
% ==================================================
\section{Functorial Semantics: Expanded Constructions and Proofs}
% ==================================================

This appendix lifts the operational semantics developed in Appendix~B to a
categorical level. The aim is to demonstrate that constraint-first meaning is not
merely an execution model but a structural semantics with well-defined invariants,
equivalences, and transformation laws. The categorical perspective clarifies which
aspects of meaning are preserved under reinterpretation and why certain semantic
failures correspond to formal type errors.

% ==================================================
\subsection{The Category of Histories}
% ==================================================

\begin{definition}[Category $\mathbf{Hist}$]
Define $\mathbf{Hist}$ as the category whose objects are event histories
$h \in \mathcal{H}$ and whose morphisms are given by prefix inclusion: there
exists a unique morphism $h \to h'$ if and only if $h \preceq h'$.
\end{definition}

\begin{lemma}[Thinness]
The category $\mathbf{Hist}$ is thin: between any two objects there exists at most
one morphism.
\end{lemma}

\begin{proof}
Prefix order is antisymmetric and transitive. Thus for any $h,h' \in \mathcal{H}$,
the relation $h \preceq h'$ uniquely determines the existence of a morphism, and
no parallel morphisms are possible.
\end{proof}

Thinness reflects the irreversibility of time: histories extend but do not branch
backward.

% ==================================================
\subsection{The Category of Constraint Spaces}
% ==================================================

\begin{definition}[Constraint Space]
A constraint space is a subset $C \subseteq \mathcal{H}$ such that whenever
$h \in C$ and $h \preceq h'$, then $h' \in C$.
\end{definition}

Constraint spaces are upward closed under prefix order and represent sets of
admissible futures.

\begin{definition}[Category $\mathbf{Constr}$]
Define $\mathbf{Constr}$ as the category whose objects are constraint spaces and
whose morphisms are monotone functions preserving inclusion.
\end{definition}

\begin{lemma}
The category $\mathbf{Constr}$ forms a complete lattice under intersection and
union.
\end{lemma}

\begin{proof}
Upward closure is preserved under arbitrary intersections and unions. Hence the
set of constraint spaces is closed under these operations, yielding a complete
lattice structure.
\end{proof}

This lattice structure reflects the combinatorial nature of constraint
composition.

% ==================================================
\subsection{The Meaning Functor}
% ==================================================

\begin{definition}[Meaning Functor]
Define a functor
\[
\mathcal{M} : \mathbf{Hist} \to \mathbf{Constr}
\]
by
\[
\mathcal{M}(h) = \mathsf{Fut}(h),
\]
and for a morphism $h \to h'$,
\[
\mathcal{M}(h \to h')(C) = C \cap \mathsf{Fut}(h').
\]
\end{definition}

\begin{theorem}[Functoriality]
The mapping $\mathcal{M}$ is a well-defined functor.
\end{theorem}

\begin{proof}
Identity: for any history $h$,
\[
\mathcal{M}(\mathrm{id}_h)(C) = C \cap \mathsf{Fut}(h) = C,
\]
since $C \subseteq \mathsf{Fut}(h)$ by definition.

Composition: if $h \preceq h' \preceq h''$, then
\[
\mathcal{M}(h \to h'') = \mathcal{M}(h' \to h'') \circ \mathcal{M}(h \to h'),
\]
because intersection with $\mathsf{Fut}(h'')$ factors through intersection with
$\mathsf{Fut}(h')$.
\end{proof}

The meaning functor assigns to each history its space of admissible continuations.

% ==================================================
\subsection{Operators as Endofunctors}
% ==================================================

\begin{definition}[Operator Endofunctor]
Each operator $\sigma$ induces an endofunctor
\[
F_\sigma : \mathbf{Constr} \to \mathbf{Constr}
\]
defined by
\[
F_\sigma(C) = C \cap \sigma(h),
\]
where $h$ is any history generating the constraint space $C$.
\end{definition}

\begin{lemma}[Well-Definedness]
The definition of $F_\sigma$ does not depend on the choice of representative
history $h$.
\end{lemma}

\begin{proof}
If $\mathsf{Fut}(h) = \mathsf{Fut}(h')$, then $h$ and $h'$ are equivalent under
prefix closure. By operator monotonicity, $\sigma(h) = \sigma(h')$, so the
intersection defining $F_\sigma(C)$ is invariant.
\end{proof}

Thus operators act uniformly on constraint spaces.

% ==================================================
\subsection{Composition and Identity}
% ==================================================

\begin{proposition}
Operator composition corresponds to functor composition:
\[
F_{\sigma_1 \circ \sigma_2} = F_{\sigma_1} \circ F_{\sigma_2}.
\]
\end{proposition}

\begin{proof}
Both sides evaluate to intersection with
$\sigma_1(h) \cap \sigma_2(h)$ for any representative history $h$.
\end{proof}

The identity operator induces the identity endofunctor on $\mathbf{Constr}$.

% ==================================================
\subsection{Reinterpretation as Natural Transformation}
% ==================================================

\begin{definition}[Reinterpretation]
A reinterpretation of a symbol is a natural transformation
\[
\eta : F_\sigma \Rightarrow F_{\sigma'}.
\]
\end{definition}

\begin{lemma}[Naturality Condition]
For all morphisms $f : C \to C'$ in $\mathbf{Constr}$,
\[
\eta_{C'} \circ F_\sigma(f) = F_{\sigma'}(f) \circ \eta_C.
\]
\end{lemma}

Failure of this condition corresponds to semantic drift: reinterpretation that
depends on context rather than preserving structural meaning.

% ==================================================
\subsection{Identity and Fixed Points}
% ==================================================

\begin{definition}[Stable Identity]
A constraint space $C$ is a stable identity if and only if
\[
F_\sigma(C) = C.
\]
\end{definition}

\begin{proposition}
Stable identities are precisely the fixed points of the operator endofunctor
system.
\end{proposition}

\begin{proof}
This follows immediately from the definition of fixed points in
$\mathbf{Constr}$.
\end{proof}

Identities are thus not primitive objects but stabilized regions of constraint
space.

% ==================================================
\subsection{Reification as a Categorical Error}
% ==================================================

\begin{definition}
Reification is the attempt to treat an endofunctor or a natural transformation as
an object of $\mathbf{Hist}$.
\end{definition}

\begin{lemma}
Such an attempt destroys functorial structure.
\end{lemma}

\begin{proof}
Objects of $\mathbf{Hist}$ are histories, not transformations. Treating
endofunctors or natural transformations as objects eliminates the distinction
between objects and morphisms, collapsing composition and violating naturality.
\end{proof}

This formalizes reification as a categorical type error.

% ==================================================
\subsection{Appendix Summary}
% ==================================================

This appendix has established the categorical structure underlying the framework,
including the categories of event histories and constraint spaces and the functor
that maps histories to their admissible futures. Within this setting, symbolic
operators were shown to correspond to endofunctors on constraint spaces, while
reinterpretation was formalized as natural transformation. Identity was
characterized as a fixed point of constraint composition, and reification was
identified as a categorical type error arising from the misclassification of
morphism-level structure as object-level content.

Taken together with the operational results developed earlier, these constructions
demonstrate that constraint-first meaning admits both an executable semantics and
an abstract categorical characterization. Crucially, the same structural
invariants are preserved across these levels of description, ensuring that the
formalism remains coherent under both execution and abstraction.

\newpage
% ==================================================
\section{Notation Glossary and Meta-Theoretic Remarks}
% ==================================================

This appendix consolidates the notation used throughout the paper and records
meta-theoretic properties of the formal system. Its purpose is expository and
clarificatory rather than developmental. No new formal machinery is introduced;
instead, this appendix makes explicit the commitments, limits, and interpretive
status of the framework as a whole.

% ==================================================
\subsection{Notation Glossary}
% ==================================================

\begin{center}
\begin{tabular}{ll}
\textbf{Symbol} & \textbf{Meaning} \\
\midrule
$\mathcal{H}$ & Set of event histories \\
$h, h'$ & Individual event histories \\
$e$ & An atomic irreversible event \\
$\prec$ & Precedence relation on events \\
$\preceq$ & Prefix order on histories \\
$\mathsf{Fut}(h)$ & Admissible futures of history $h$ \\
$\mathcal{O}$ & Semantic type of operators \\
$\mathcal{R}$ & Semantic type of referential terms \\
$\sigma, \tau$ & Operators (future-constraint functions) \\
$\Sigma$ & Finite multiset of active operators \\
$(h,\Sigma)$ & Operational configuration \\
$\mathbf{Hist}$ & Category of event histories \\
$\mathbf{Constr}$ & Category of constraint spaces \\
$\mathcal{M}$ & Meaning functor $\mathbf{Hist} \to \mathbf{Constr}$ \\
$F_\sigma$ & Endofunctor induced by operator $\sigma$ \\
$\eta$ & Natural transformation (reinterpretation) \\
$\circ$ & Operator or functor composition \\
$\varnothing$ & Empty future space (semantic deadlock)
\end{tabular}
\end{center}

\begin{remark}
All symbols are typed implicitly by context. Ambiguity of semantic type is treated
as a formal error rather than an interpretive ambiguity.
\end{remark}

% ==================================================
\subsection{Meta-Theoretic Status of the System}
% ==================================================

This subsection records what has been established formally and what remains
outside the scope of proof.

\paragraph{Consistency.}
The system is consistent relative to standard set theory. No rule introduces
contradiction because no rule asserts propositions. All semantics is structural,
monotone, and eliminative rather than generative.

\paragraph{Soundness and Completeness.}
Soundness and completeness were proven with respect to the operational semantics
in Appendix~B. These results are internal: they concern faithful execution of
constraint restriction rather than correspondence to an external model of truth.

\paragraph{Decidability.}
Decidability of non-empty future space depends on the class of operators admitted.
For finite operator sets with computable constraint predicates, deadlock detection
is decidable. In the general case, undecidability is expected.

\paragraph{Expressive Power.}
The system is strictly more expressive than propositional semantics, but strictly
less expressive than unrestricted higher-order logic. Its expressive strength
derives from irreversibility and constraint composition rather than from
quantification over truth values.

% ==================================================
\subsection{Negative Results and Non-Claims}
% ==================================================

For the avoidance of misunderstanding, it is important to state explicitly what
the framework does not claim. It does not assert that all disagreement is semantic
in origin, nor does it suggest that referential semantics can be eliminated or
replaced wholesale. It does not maintain that operator-level interpretation is
sufficient to resolve every form of conflict, and it does not propose that formal
models of constraint provide an exhaustive account of lived experience, social
dynamics, or psychological motivation.

These negative claims are not incidental qualifications. They are structural
boundaries that prevent theoretical overreach and clarify the specific explanatory
target of the framework. By delineating what lies outside its scope, the analysis
preserves both its neutrality and its analytical force.

\begin{remark}
The framework diagnoses a specific semantic failure mode. It is not a universal
solvent.
\end{remark}

% ==================================================
\subsection{Formalism and Interpretation}
% ==================================================

The formal system constrains admissible interpretations but does not determine
them uniquely. Multiple narratives, institutions, or practices may instantiate
the same constraint topology.

This explains how semantic equivalence can coexist with surface disagreement.

\begin{proposition}
If two symbolic systems induce naturally isomorphic meaning functors, then they
are semantically equivalent with respect to admissible futures.
\end{proposition}

\begin{proof}
Natural isomorphism preserves all structural relations in $\mathbf{Constr}$,
including inclusion, intersection, and fixed points. Hence the induced spaces of
admissible futures coincide.
\end{proof}

% ==================================================
\subsection{Reification at the Meta-Level}
% ==================================================

Reification can itself be understood meta-theoretically as a collapse of typing
discipline.

\begin{definition}[Meta-Reification]
Meta-reification occurs when the semantic framework itself is treated as a
descriptive ontology rather than as a constraint model.
\end{definition}

This appendix explicitly avoids that error. The formalism does not describe what
exists; it specifies how meaning constrains continuation.

% ==================================================
\subsection{Appendix Summary}
% ==================================================

This appendix has served to consolidate the formal presentation of the framework
by fixing notation unambiguously and recording the meta-theoretic status of the
system. It has clarified which structural guarantees are established by the
formalism and which questions are intentionally left open, thereby preventing
overextension or misinterpretation of the results. In doing so, it has made
explicit the scope conditions under which the framework is intended to apply and
the limits beyond which its claims do not extend.

In addition, the appendix has articulated criteria for semantic equivalence in
terms of preserved constraint structure, rather than surface representation or
interpretive narrative. Finally, it has introduced a principled safeguard against
reification at the level of the theory itself, ensuring that the formal apparatus
is understood as a model of constraint and coordination rather than as a
descriptive ontology.

Taken together, these clarifications complete the formal specification of the
constraint-first semantic framework and establish the conditions under which its
results should be interpreted.

\newpage
\bibliographystyle{plainnat}
\bibliography{constraint-before-reference}



\end{document}
